\documentclass[11pt,a4paper]{article}

\usepackage[utf8]{inputenc}
\usepackage[indonesian]{babel}
\usepackage{amsmath,amssymb,amsfonts,amsthm}
\usepackage{geometry}
\usepackage{tcolorbox}
\usepackage{booktabs}
\usepackage{multirow}
\usepackage{adjustbox}
\usepackage{enumitem}
\usepackage{hyperref}
\usepackage{tikz}
\usepackage{pgfplots}
\pgfplotsset{compat=1.18}

\geometry{
    top=2.5cm,
    bottom=2.5cm,
    left=2.5cm,
    right=2.5cm
}

\tcbuselibrary{skins,breakable}

% Colors
\definecolor{primaryblue}{RGB}{24, 76, 120}
\definecolor{accentblue}{RGB}{70, 130, 180}
\definecolor{boxbg}{RGB}{245, 248, 252}
\definecolor{darkgreen}{RGB}{34, 139, 34}
\definecolor{darkred}{RGB}{178, 34, 34}
\definecolor{darkorange}{RGB}{205, 92, 0}

% Custom theorem/box styles
\newtcolorbox{definitionbox}[1][]{
    enhanced,
    breakable,
    colback=boxbg,
    colframe=primaryblue,
    fonttitle=\bfseries\sffamily,
    title={#1},
    boxrule=0.8pt,
    titlerule=0pt,
    coltitle=white,
    colbacktitle=primaryblue,
    arc=3mm,
    boxsep=2mm
}

\newtcolorbox{examplebox}[1][]{
    enhanced,
    breakable,
    colback=white,
    colframe=accentblue,
    fonttitle=\bfseries\sffamily,
    title={#1},
    boxrule=0.8pt,
    arc=2mm,
    coltitle=white,
    colbacktitle=accentblue,
    boxsep=2mm
}

\newtcolorbox{notebox}[1][]{
    enhanced,
    breakable,
    colback=white,
    colframe=gray!60,
    fonttitle=\bfseries\sffamily,
    title={#1},
    boxrule=0.5pt,
    arc=2mm,
    coltitle=black!80,
    colbacktitle=gray!20,
    boxsep=2mm
}

\hypersetup{
    colorlinks=true,
    linkcolor=primaryblue,
    urlcolor=accentblue,
    citecolor=primaryblue
}

\title{\textbf{\LARGE Catatan Kuliah Teori Suku Bunga}\\[0.3em]
\Large Pertemuan 07: Anuitas yang Lebih Umum (Bagian 1)\\
\large (\textit{General Annuities: Payment Frequency Differs from Interest Conversion Frequency \& Continuous Annuities})}
\author{\textbf{Referensi:} Stephen G. Kellison, \textit{The Theory of Interest} (3rd ed., Bab 4) \& Slide Perkuliahan ITB}
\date{\today}

\begin{document}

\maketitle
\tableofcontents
\vspace{1em}
\noindent\rule{\textwidth}{0.4pt}
\vspace{1.5em}

\section{Pendahuluan \& Konsep Dasar Tiga Periode Waktu}

Pada Bab 05 dan 06, kita telah mempelajari \textbf{Anuitas Sederhana (\textit{Simple Annuities})}, di mana terdapat asumsi dasar yang sangat penting:
\begin{center}
    \textbf{Frekuensi pembayaran cicilan/setoran persis sama dengan frekuensi konversi (pemajemukan) suku bunga.}
\end{center}
Sebagai contoh, jika bunga dimajemukkan secara tahunan, maka pembayaran cicilan juga dilakukan setahun sekali; jika bunga dimajemukkan secara bulanan, pembayaran cicilan juga bulanan.

Namun dalam praktik dunia keuangan, perbankan, dan aktuarial modern, kondisi ideal ini sering kali tidak terpenuhi. Sebagai contoh:
\begin{itemize}[leftmargin=1.5em]
    \item Seseorang membayar premi asuransi atau menabung cicilan \textbf{setiap bulan}, tetapi bank memajemukkan bunga secara \textbf{semesteran} atau \textbf{tahunan}.
    \item Suatu perusahaan menerima kupon obligasi \textbf{setiap semester}, tetapi dana tersebut diinvestasikan kembali pada rekening dengan pemajemukan bunga \textbf{bulanan}.
    \item Aliran kas bisnis atau penerimaan premi terjadi terus-menerus sepanjang waktu (\textbf{secara kontinu}).
\end{itemize}

Oleh karena itu, kita memerlukan teori \textbf{Anuitas yang Lebih Umum (\textit{General Annuities})}.

\begin{definitionbox}[Tiga Dimensi Periode Waktu dalam Transaksi Anuitas]
Dalam menganalisis anuitas umum, kita wajib mengidentifikasi tiga periode waktu yang berbeda:
\begin{enumerate}[label=\textbf{\arabic*.}]
    \item \textbf{Periode Konversi Suku Bunga (\textit{Interest Conversion Period}):} Interval waktu di mana bunga dihitung dan ditambahkan ke pokok (\textit{compounded}), misalnya tahunan, semesteran, kuartalan, atau bulanan.
    \item \textbf{Periode Pembayaran Dana/Cicilan (\textit{Payment Period}):} Interval waktu antara dua pembayaran berturut-turut, misalnya setiap 3 bulan, setiap bulan, atau setiap minggu.
    \item \textbf{Periode / Jangka Waktu Investasi (\textit{Investment Term}):} Total rentang waktu keseluruhan transaksi anuitas berlangsung, misalnya 4 tahun, 10 tahun, atau selamanya ($\infty$).
\end{enumerate}
\end{definitionbox}

\subsection{Dua Pendekatan Penyelesaian Masalah}

Ketika periode pembayaran dan periode konversi suku bunga \textbf{tidak sama}, terdapat dua pendekatan utama dalam teori suku bunga:

\begin{center}
\begin{tikzpicture}[scale=0.95, >=stealth]
    \node[draw=primaryblue, fill=boxbg, rounded corners, thick, text width=6.5cm, align=center, inner sep=8pt] (p1) at (-4,0) {
        \textbf{\large Pendekatan 1: Perhitungan Numerik}\\[0.3em]
        \small (\textit{Effective Interest Rate Equivalence})\\[0.5em]
        Mengonversi suku bunga menjadi suku bunga efektif per periode pembayaran ($j$), lalu menggunakan rumus anuitas standar.
    };
    
    \node[draw=accentblue, fill=white, rounded corners, thick, text width=6.5cm, align=center, inner sep=8pt] (p2) at (4,0) {
        \textbf{\large Pendekatan 2: Analisis Aljabar}\\[0.3em]
        \small (\textit{Actuarial Formulaic Approach})\\[0.5em]
        Mempertahankan periode konversi bunga sebagai unit waktu acuan dan membangun rumus analitik baru berbasis simbol aktuaria ($a_{\overline{n}|}/s_{\overline{k}|}$, $a_{\overline{n}|}^{(m)}$, $\bar{a}_{\overline{n}|}$).
    };
\end{tikzpicture}
\end{center}

---

\section{Pendekatan 1: Perhitungan Numerik (Suku Bunga Ekuivalen)}

\subsection{Prinsip Dasar Ekuivalensi Suku Bunga}
Prinsip fundamental dari pendekatan ini adalah:
\begin{quote}
    \textit{"Ubah satuan waktu masalah agar sesuai dengan interval pembayaran, lalu cari suku bunga efektif per periode pembayaran tersebut."}
\end{quote}

Jika diketahui suku bunga efektif per periode konversi bunga adalah $i$, dan dalam satu periode pembayaran terdapat kelipatan atau fraksi periode bunga, kita mencari suku bunga efektif per periode pembayaran, dinotasikan dengan $j$, menggunakan persamaan akumulasi:
\begin{equation}
    (1 + j) = (1 + i)^t
\end{equation}
di mana $t$ adalah lamanya satu periode pembayaran diukur dalam satuan periode konversi bunga.

Setelah suku bunga $j$ diperoleh, seluruh perhitungan nilai sekarang ($PV$) dan nilai terakumulasi ($AV/FV$) dapat langsung diselesaikan menggunakan rumus anuitas dasar berjangka $N$ kali pembayaran:
\begin{equation}
    PV = R \cdot a_{\overline{N}|j} = R \left( \frac{1 - (1+j)^{-N}}{j} \right)
\end{equation}
\begin{equation}
    FV = R \cdot s_{\overline{N}|j} = R \left( \frac{(1+j)^N - 1}{j} \right)
\end{equation}
Untuk anuitas di muka (\textit{annuity-due}), kalikan dengan faktor $(1+j)$ atau gunakan penyebut tingkat diskon $d_j = \frac{j}{1+j}$.

---

\subsection{Contoh Soal 1: Nilai Akumulasi Tabungan Bertahap}

\begin{examplebox}[Contoh 1 (Slide Perkuliahan Hal. 3--5)]
Cari nilai akumulasi dana investasi pada akhir tahun ke-4 di mana:
\begin{itemize}
    \item $\$100$ disimpan pada \textbf{awal dari setiap kuartal} untuk 2 tahun pertama.
    \item $\$200$ disimpan pada \textbf{awal dari setiap kuartal} untuk 2 tahun berikutnya.
    \item Tingkat suku bunga nominal yang berlaku adalah $12\%$ \textit{convertible monthly} ($i^{(12)} = 12\%$).
\end{itemize}
\end{examplebox}

\begin{proof}[\textbf{Penyelesaian dengan Pendekatan Numerik:}]
\textbf{Langkah 1: Identifikasi Data dan Satuan Waktu}
\begin{itemize}
    \item Total periode investasi: $4\text{ tahun}$.
    \item Periode konversi bunga: Bulanan ($m=12$), sehingga suku bunga efektif per bulan adalah:
    \begin{equation*}
        i_{\text{bulan}} = \frac{i^{(12)}}{12} = \frac{12\%}{12} = 1\% = 0.01
    \end{equation*}
    \item Periode pembayaran: Kuartalan (setiap 3 bulan).
    \item Banyaknya pembayaran total: $(4 \times 2) + (4 \times 2) = 8 + 8 = 16\text{ kali pembayaran}$.
\end{itemize}

\textbf{Langkah 2: Menghitung Suku Bunga Efektif per Kuartal ($j$)}
Karena 1 kuartal terdiri dari 3 bulan, maka faktor akumulasi 1 kuartal harus sama dengan akumulasi 3 bulan:
\begin{equation*}
    1 + j = (1 + i_{\text{bulan}})^3 = (1 + 0.01)^3 = 1.030301
\end{equation*}
\begin{equation*}
    j = 1.030301 - 1 = 0.030301 \quad (3.0301\%\text{ per kuartal})
\end{equation*}
Tingkat diskon efektif per kuartal ($d_j$) adalah:
\begin{equation*}
    d_j = \frac{j}{1+j} = \frac{0.030301}{1.030301} \approx 0.0294099
\end{equation*}

\textbf{Langkah 3: Dekomposisi Arus Kas}
Penyetoran dilakukan di awal kuartal (\textit{annuity-due}). Kita dapat membagi aliran kas menjadi dua lapis anuitas:
\begin{enumerate}
    \item Anuitas A: Setoran $\$100$ per kuartal selama 4 tahun penuh (16 kuartal).
    \item Anuitas B: Setoran tambahan $\$100$ per kuartal selama 2 tahun terakhir (8 kuartal).
\end{enumerate}

\begin{center}
\begin{tikzpicture}[scale=0.9, >=stealth]
    \draw[->, thick] (0,0) -- (13,0) node[right] {Kuartal};
    \foreach \x in {0,1,2,3,4,5,6,7,8,9,10,11,12} {
        \draw[thick] (\x*0.9,0.1) -- (\x*0.9,-0.1);
    }
    \node[below] at (0, -0.15) {\small 0};
    \node[below] at (7.2, -0.15) {\small 8 (Thn 2)};
    \node[below] at (14.4*0.8, -0.15) {\small 16 (Thn 4)};
    
    % Layer 1: $100 for 16 quarters
    \foreach \x in {0,1.8,3.6,5.4,7.2,9.0,10.8} {
        \draw[->, thick, primaryblue] (\x, 0.2) -- (\x, 0.9) node[above] {\small $\$100$};
    }
    % Layer 2: $100 for last 8 quarters
    \foreach \x in {7.2,9.0,10.8} {
        \draw[->, thick, darkred] (\x, 1.2) -- (\x, 1.9) node[above] {\small $+\$100$};
    }
\end{tikzpicture}
\end{center}

\textbf{Langkah 4: Menghitung Nilai Akumulasi pada Akhir Tahun ke-4 ($t = 16$ kuartal)}
\begin{align*}
    FV &= 100 \cdot \ddot{s}_{\overline{16}|j} + 100 \cdot \ddot{s}_{\overline{8}|j} \\
       &= 100 \left[ \frac{(1+j)^{16} - 1}{d_j} + \frac{(1+j)^8 - 1}{d_j} \right] \\
       &= 100 \left[ \frac{(1.030301)^{16} - 1}{\frac{0.030301}{1.030301}} + \frac{(1.030301)^8 - 1}{\frac{0.030301}{1.030301}} \right] \\
       &= 100 \left[ \frac{1.6122261 - 1}{0.0294099} + \frac{1.2697346 - 1}{0.0294099} \right] \\
       &= 100 \left[ \frac{0.6122261}{0.0294099} + \frac{0.2697346}{0.0294099} \right] \\
       &= 100 \left[ 20.81704 + 9.17157 \right] = 100 (29.98861) = \mathbf{2998.86149}
\end{align*}
Jadi, nilai akumulasi dana investasi tersebut adalah \textbf{\$2.998,86}.
\end{proof}

---

\subsection{Contoh Soal 2: Menghitung Besarnya Angsuran Pinjaman}

\begin{examplebox}[Contoh 2 (Slide Perkuliahan Hal. 6--7)]
Suatu pinjaman sebesar $\$3.000$ dilunasi dengan angsuran berkala pada setiap \textbf{akhir kuartal} selama 5 tahun. Jika suku bunga pinjaman adalah $10\%$ \textit{convertible semiannually} ($i^{(2)} = 10\%$), hitunglah besarnya angsuran per kuartal ($R$).
\end{examplebox}

\begin{proof}[\textbf{Penyelesaian dengan Pendekatan Numerik:}]
\textbf{Langkah 1: Parameter Transaksi}
\begin{itemize}
    \item Pokok pinjaman ($PV$): $\$3.000$.
    \item Jangka waktu: 5 tahun.
    \item Frekuensi pembayaran: Kuartalan $\implies$ Total pembayaran $N = 5 \times 4 = 20\text{ kali}$.
    \item Suku bunga konversi semesteran: $i_{\text{semester}} = \frac{i^{(2)}}{2} = \frac{10\%}{2} = 5\% = 0.05$.
\end{itemize}

\textbf{Langkah 2: Menghitung Suku Bunga Efektif per Kuartal ($j$)}
Karena dalam 1 semester terdapat 2 kuartal:
\begin{equation*}
    (1 + j)^2 = 1 + i_{\text{semester}} = 1.05 \implies 1 + j = (1.05)^{1/2}
\end{equation*}
\begin{equation*}
    j = (1.05)^{1/2} - 1 = 1.024695076 - 1 \approx 0.024695 \quad (2.4695\%\text{ per kuartal})
\end{equation*}

\textbf{Langkah 3: Persamaan Nilai Sekarang (\textit{Equation of Value})}
Karena pembayaran dilakukan di akhir kuartal, berlaku formula \textit{annuity-immediate}:
\begin{equation*}
    PV = R \cdot a_{\overline{20}|j} \implies 3000 = R \cdot \left( \frac{1 - (1+j)^{-20}}{j} \right)
\end{equation*}
Hitung faktor diskonto:
\begin{equation*}
    (1+j)^{-20} = (1.05)^{-10} \approx 0.61391325
\end{equation*}
\begin{equation*}
    a_{\overline{20}|j} = \frac{1 - 0.61391325}{0.024695076} = \frac{0.38608675}{0.024695076} \approx 15.63416
\end{equation*}
Maka besarnya angsuran per kuartal adalah:
\begin{equation*}
    R = \frac{3000}{a_{\overline{20}|j}} = \frac{3000}{15.63416} = \mathbf{191.8874} \approx \mathbf{\$191.89}
\end{equation*}
\end{proof}

---

\section{Kasus 1 Aljabar: Pembayaran Lebih Jarang daripada Pemajemukan Bunga (\texorpdfstring{$k > 1$}{k > 1})}

\subsection{Definisi dan Notasi Parameter}
Pada kasus ini, interval pembayaran lebih panjang daripada interval pemajemukan bunga. Misalkan:
\begin{itemize}[leftmargin=1.5em]
    \item $k$: Banyaknya periode konversi suku bunga dalam satu kali pembayaran ($k \in \mathbb{Z}^+, k > 1$).
    \item $n$: Total jangka waktu anuitas yang diukur dalam \textbf{satuan periode konversi suku bunga}.
    \item $m$: Total banyaknya pembayaran yang dilakukan sepanjang masa anuitas, di mana:
    \begin{equation}
        m = \frac{n}{k} \in \mathbb{Z}^+
    \end{equation}
    \item $i$: Suku bunga efektif per periode konversi suku bunga.
    \item $v = (1+i)^{-1}$: Faktor diskonto per periode konversi suku bunga.
\end{itemize}

---

\subsection{Anuitas Akhir (\textit{Annuity-Immediate}) untuk \texorpdfstring{$k > 1$}{k > 1}}

Misalkan pembayaran sebesar $1$ dilakukan pada setiap akhir $k$ periode bunga, yaitu pada titik waktu:
\begin{equation*}
    t = k, 2k, 3k, \dots, mk = n
\end{equation*}

\begin{center}
\begin{tikzpicture}[scale=1.0, >=stealth]
    % Time axis
    \draw[->, thick] (-0.5,0) -- (11.5,0) node[right] {Periode Bunga ($t$)};
    
    % Ticks and sub-ticks
    \draw[thick] (0,0.15) -- (0,-0.15) node[below] {\small $0$};
    \draw[thick] (1,0.08) -- (1,-0.08) node[below] {\tiny $1$};
    \draw[thick] (2,0.08) -- (2,-0.08) node[below] {\tiny $2$};
    \draw[dotted, thick] (2.2,0) -- (2.8,0);
    \draw[thick] (3,0.15) -- (3,-0.15) node[below] {\small $k$};
    
    \draw[thick] (4,0.08) -- (4,-0.08);
    \draw[dotted, thick] (4.2,0) -- (5.8,0);
    \draw[thick] (6,0.15) -- (6,-0.15) node[below] {\small $2k$};
    
    \draw[dotted, thick] (6.5,0) -- (9.5,0);
    \draw[thick] (10,0.15) -- (10,-0.15) node[below] {\small $mk=n$};
    
    % Cash flows at k, 2k, ..., n
    \draw[->, thick, primaryblue] (3, 0.2) -- (3, 1.2) node[above] {\textbf{\$1}};
    \draw[->, thick, primaryblue] (6, 0.2) -- (6, 1.2) node[above] {\textbf{\$1}};
    \draw[->, thick, primaryblue] (10, 0.2) -- (10, 1.2) node[above] {\textbf{\$1}};
    
    % PV & FV
    \draw[<-, thick, darkred] (0, -0.6) -- (0, -1.2) node[below] {\small $PV = \frac{a_{\overline{n}|i}}{s_{\overline{k}|i}}$};
    \draw[<-, thick, darkgreen] (10, -0.6) -- (10, -1.2) node[below] {\small $FV = \frac{s_{\overline{n}|i}}{s_{\overline{k}|i}}$};
\end{tikzpicture}
\end{center}

\subsubsection{Penurunan Matematis Nilai Sekarang (\textit{Present Value})}
Nilai sekarang pada $t = 0$ adalah jumlahan deret diskonto:
\begin{equation}
    PV = v^k + v^{2k} + v^{3k} + \dots + v^{mk} = v^k + v^{2k} + \dots + v^n
\end{equation}
Deret di atas merupakan deret geometri dengan suku pertama $a_1 = v^k$, rasio $r = v^k$, dan banyak suku $m$:
\begin{equation}
    PV = v^k \left( \frac{1 - (v^k)^m}{1 - v^k} \right) = v^k \left( \frac{1 - v^n}{1 - v^k} \right)
\end{equation}
Kalikan pembilang dan penyebut dengan $v^{-k} = (1+i)^k$:
\begin{equation}
    PV = \frac{1 - v^n}{(1+i)^k - 1}
\end{equation}
Bagi pembilang dan penyebut dengan suku bunga $i$:
\begin{equation}
    PV = \frac{\frac{1 - v^n}{i}}{\frac{(1+i)^k - 1}{i}} = \mathbf{\frac{a_{\overline{n}|i}}{s_{\overline{k}|i}}}
\end{equation}

\subsubsection{Penurunan Nilai Terakumulasi (\textit{Accumulated Value})}
Nilai akumulasi pada titik $t = n$ diperoleh dengan mengalikan $PV$ dengan $(1+i)^n$:
\begin{equation}
    FV = PV \cdot (1+i)^n = \frac{a_{\overline{n}|i} \cdot (1+i)^n}{s_{\overline{k}|i}} = \mathbf{\frac{s_{\overline{n}|i}}{s_{\overline{k}|i}}}
\end{equation}

\begin{definitionbox}[Rumus Pokok: Pembayaran 1 Setiap Akhir $k$ Periode Bunga]
\begin{equation}
    PV = \frac{a_{\overline{n}|i}}{s_{\overline{k}|i}} \qquad \text{dan} \qquad FV = \frac{s_{\overline{n}|i}}{s_{\overline{k}|i}}
\end{equation}
\textbf{Interpretasi Aktuarial yang Sangat Cantik:}\\
Pembayaran sebesar $1$ pada akhir setiap $k$ periode bunga \textbf{ekuivalen secara nilai waktu} dengan serangkaian pembayaran sebesar $\frac{1}{s_{\overline{k}|i}}$ pada setiap akhir periode bunga tunggal selama $n$ periode. Oleh karena itu, $PV = \frac{1}{s_{\overline{k}|i}} \cdot a_{\overline{n}|i} = \frac{a_{\overline{n}|i}}{s_{\overline{k}|i}}$.
\end{definitionbox}

---

\subsection{Anuitas Awal (\textit{Annuity-Due}) untuk \texorpdfstring{$k > 1$}{k > 1}}

Misalkan pembayaran sebesar $1$ dilakukan pada setiap awal $k$ periode bunga, yaitu pada titik waktu:
\begin{equation*}
    t = 0, k, 2k, \dots, n-k
\end{equation*}

\begin{center}
\begin{tikzpicture}[scale=1.0, >=stealth]
    \draw[->, thick] (-0.5,0) -- (11.5,0) node[right] {Periode Bunga ($t$)};
    
    \draw[thick] (0,0.15) -- (0,-0.15) node[below] {\small $0$};
    \draw[thick] (3,0.15) -- (3,-0.15) node[below] {\small $k$};
    \draw[thick] (6,0.15) -- (6,-0.15) node[below] {\small $2k$};
    \draw[thick] (8.5,0.15) -- (8.5,-0.15) node[below] {\small $n-k$};
    \draw[thick] (10.5,0.15) -- (10.5,-0.15) node[below] {\small $n$};
    
    % Cash flows at 0, k, 2k, ..., n-k
    \draw[->, thick, primaryblue] (0, 0.2) -- (0, 1.2) node[above] {\textbf{\$1}};
    \draw[->, thick, primaryblue] (3, 0.2) -- (3, 1.2) node[above] {\textbf{\$1}};
    \draw[->, thick, primaryblue] (6, 0.2) -- (6, 1.2) node[above] {\textbf{\$1}};
    \draw[->, thick, primaryblue] (8.5, 0.2) -- (8.5, 1.2) node[above] {\textbf{\$1}};
    
    % PV & FV
    \draw[<-, thick, darkred] (0, -0.6) -- (0, -1.2) node[below] {\small $PV = \frac{a_{\overline{n}|i}}{a_{\overline{k}|i}}$};
    \draw[<-, thick, darkgreen] (10.5, -0.6) -- (10.5, -1.2) node[below] {\small $FV = \frac{s_{\overline{n}|i}}{a_{\overline{k}|i}}$};
\end{tikzpicture}
\end{center}

\subsubsection{Penurunan Matematis Nilai Sekarang}
Nilai sekarang pada $t = 0$:
\begin{equation}
    PV = 1 + v^k + v^{2k} + \dots + v^{n-k} = \frac{1 - (v^k)^m}{1 - v^k} = \frac{1 - v^n}{1 - v^k}
\end{equation}
Bagi pembilang dan penyebut dengan $i$:
\begin{equation}
    PV = \frac{\frac{1 - v^n}{i}}{\frac{1 - v^k}{i}} = \mathbf{\frac{a_{\overline{n}|i}}{a_{\overline{k}|i}}}
\end{equation}

\subsubsection{Penurunan Nilai Terakumulasi}
Akumulasi pada $t = n$:
\begin{equation}
    FV = PV \cdot (1+i)^n = \frac{a_{\overline{n}|i}(1+i)^n}{a_{\overline{k}|i}} = \mathbf{\frac{s_{\overline{n}|i}}{a_{\overline{k}|i}}}
\end{equation}

\begin{definitionbox}[Rumus Pokok: Pembayaran 1 Setiap Awal $k$ Periode Bunga]
\begin{equation}
    PV = \frac{a_{\overline{n}|i}}{a_{\overline{k}|i}} \qquad \text{dan} \qquad FV = \frac{s_{\overline{n}|i}}{a_{\overline{k}|i}}
\end{equation}
\textbf{Interpretasi Aktuarial:}\\
Pembayaran $1$ pada awal setiap interval sepanjang $k$ periode bunga dapat digantikan oleh pembayaran berkala sebesar $\frac{1}{a_{\overline{k}|i}}$ pada setiap akhir periode bunga tunggal. Oleh karena itu, $PV = \frac{1}{a_{\overline{k}|i}} \cdot a_{\overline{n}|i} = \frac{a_{\overline{n}|i}}{a_{\overline{k}|i}}$.
\end{definitionbox}

---

\subsection{Perpetuitas untuk Kasus \texorpdfstring{$k > 1$}{k > 1}}

Jika pembayaran berlangsung selamanya ($m \to \infty$ sehingga $n \to \infty$), kita peroleh nilai sekarang perpetuitas:

\begin{equation}
    PV_{\text{Perp-Immediate}} = \lim_{n \to \infty} \frac{a_{\overline{n}|i}}{s_{\overline{k}|i}} = \frac{\frac{1}{i}}{s_{\overline{k}|i}} = \mathbf{\frac{1}{i \cdot s_{\overline{k}|i}} = \frac{1}{(1+i)^k - 1}}
\end{equation}

\begin{equation}
    PV_{\text{Perp-Due}} = \lim_{n \to \infty} \frac{a_{\overline{n}|i}}{a_{\overline{k}|i}} = \frac{\frac{1}{i}}{a_{\overline{k}|i}} = \mathbf{\frac{1}{i \cdot a_{\overline{k}|i}} = \frac{1}{1 - v^k}}
\end{equation}

---

\subsection{Aplikasi: Menghitung Ulang Contoh 1 dengan Pendekatan Aljabar}

\begin{examplebox}[Penyelesaian Aljabar Contoh 1 (Slide Perkuliahan Hal. 11)]
Pada Contoh 1, suku bunga adalah $1\%$ per bulan ($i = 0.01$).
\begin{itemize}
    \item Interval pembayaran: 3 bulan ($k = 3$).
    \item Lapisan 1: Setoran $\$100$ di awal kuartal selama 4 tahun ($n = 4 \times 12 = 48\text{ bulan}$).
    \item Lapisan 2: Setoran tambahan $\$100$ di awal kuartal selama 2 tahun terakhir ($n = 2 \times 12 = 24\text{ bulan}$).
\end{itemize}
\end{examplebox}

\begin{proof}[\textbf{Langkah Perhitungan:}]
Menggunakan rumus akumulasi anuitas di muka untuk $k > 1$, yaitu $FV = \frac{s_{\overline{n}|i}}{a_{\overline{k}|i}}$:
\begin{align*}
    FV &= 100 \cdot \frac{s_{\overline{48}|0.01}}{a_{\overline{3}|0.01}} + 100 \cdot \frac{s_{\overline{24}|0.01}}{a_{\overline{3}|0.01}} \\
       &= 100 \left( \frac{s_{\overline{48}|0.01} + s_{\overline{24}|0.01}}{a_{\overline{3}|0.01}} \right)
\end{align*}
Hitung masing-masing nilai aktuaria:
\begin{align*}
    s_{\overline{48}|0.01} &= \frac{(1.01)^{48} - 1}{0.01} = \frac{1.6122261 - 1}{0.01} = 61.22261 \\
    s_{\overline{24}|0.01} &= \frac{(1.01)^{24} - 1}{0.01} = \frac{1.2697346 - 1}{0.01} = 26.97346 \\
    a_{\overline{3}|0.01}  &= \frac{1 - (1.01)^{-3}}{0.01} = \frac{1 - 0.970590}{0.01} = 2.94099
\end{align*}
Substitusikan:
\begin{equation*}
    FV = 100 \left( \frac{61.22261 + 26.97346}{2.94099} \right) = 100 \left( \frac{88.19607}{2.94099} \right) = 100 (29.98857) = \mathbf{2998.86}
\end{equation*}
Hasil ini persis identik dengan Pendekatan Numerik sebelumnya!
\end{proof}

---

\section{Kasus 2 Aljabar: Pembayaran Lebih Sering daripada Pemajemukan Bunga (\texorpdfstring{$m > 1$}{m > 1})}

\subsection{Definisi dan Notasi Standar Aktuaria}
Pada kasus ini, interval pembayaran lebih sering daripada interval konversi suku bunga.
\begin{itemize}[leftmargin=1.5em]
    \item Periode konversi suku bunga dijadikan sebagai \textbf{unit waktu acuan} ($t = 1, 2, \dots, n$).
    \item $m$: Banyaknya frekuensi pembayaran dalam satu periode konversi suku bunga ($m \in \mathbb{Z}^+, m > 1$).
    \item Total banyaknya pembayaran: $mn$ kali.
    \item Total pembayaran dalam satu periode bunga bernilai $1$, sehingga besarnya setiap kali pembayaran adalah $\mathbf{\frac{1}{m}}$.
    \item $i$: Suku bunga efektif per periode konversi bunga.
    \item $i^{(m)}$: Suku bunga nominal yang dapat dikonversi $m$ kali per periode bunga, di mana:
    \begin{equation}
        i^{(m)} = m \left[ (1+i)^{1/m} - 1 \right] \iff 1 + \frac{i^{(m)}}{m} = (1+i)^{1/m}
    \end{equation}
    \item $d^{(m)}$: Tingkat diskonto nominal yang dapat dikonversi $m$ kali per periode bunga, di mana:
    \begin{equation}
        d^{(m)} = m \left[ 1 - (1-d)^{1/m} \right] = m \left[ 1 - v^{1/m} \right] \iff 1 - \frac{d^{(m)}}{m} = v^{1/m}
    \end{equation}
\end{itemize}

---

\subsection{Anuitas Akhir: Notasi \texorpdfstring{$a_{\overline{n}|}^{(m)}$}{a\_n\textasciicircum(m)} dan \texorpdfstring{$s_{\overline{n}|}^{(m)}$}{s\_n\textasciicircum(m)}}

Misalkan pembayaran sebesar $\frac{1}{m}$ dilakukan pada setiap akhir sub-periode $\frac{1}{m}$, yaitu pada waktu:
\begin{equation*}
    t = \frac{1}{m}, \frac{2}{m}, \frac{3}{m}, \dots, n
\end{equation*}

\begin{center}
\begin{tikzpicture}[scale=1.0, >=stealth]
    \draw[->, thick] (-0.5,0) -- (11.5,0) node[right] {Periode Bunga ($t$)};
    
    \draw[thick] (0,0.15) -- (0,-0.15) node[below] {\small $0$};
    \draw[thick] (1,0.08) -- (1,-0.08) node[below] {\tiny $1/m$};
    \draw[thick] (2,0.08) -- (2,-0.08) node[below] {\tiny $2/m$};
    \draw[dotted, thick] (2.2,0) -- (2.8,0);
    \draw[thick] (3,0.15) -- (3,-0.15) node[below] {\small $1$};
    \draw[thick] (6,0.15) -- (6,-0.15) node[below] {\small $2$};
    \draw[dotted, thick] (6.5,0) -- (9.5,0);
    \draw[thick] (10,0.15) -- (10,-0.15) node[below] {\small $n$};
    
    % Cash flows of 1/m
    \draw[->, thick, primaryblue] (1, 0.2) -- (1, 1.0) node[above] {\small $\frac{1}{m}$};
    \draw[->, thick, primaryblue] (2, 0.2) -- (2, 1.0) node[above] {\small $\frac{1}{m}$};
    \draw[->, thick, primaryblue] (3, 0.2) -- (3, 1.0) node[above] {\small $\frac{1}{m}$};
    \draw[->, thick, primaryblue] (10, 0.2) -- (10, 1.0) node[above] {\small $\frac{1}{m}$};
    
    % PV & FV
    \draw[<-, thick, darkred] (0, -0.6) -- (0, -1.2) node[below] {\small $PV = a_{\overline{n}|}^{(m)}$};
    \draw[<-, thick, darkgreen] (10, -0.6) -- (10, -1.2) node[below] {\small $FV = s_{\overline{n}|}^{(m)}$};
\end{tikzpicture}
\end{center}

\subsubsection{Penurunan Matematis Nilai Sekarang \texorpdfstring{$a_{\overline{n}|}^{(m)}$}{a\_n\textasciicircum(m)}}
\begin{align}
    a_{\overline{n}|}^{(m)} &= \frac{1}{m} \left( v^{1/m} + v^{2/m} + v^{3/m} + \dots + v^{n - 1/m} + v^n \right) \nonumber\\[0.3em]
    &= \frac{1}{m} \cdot v^{1/m} \left( \frac{1 - (v^{1/m})^{mn}}{1 - v^{1/m}} \right) = \frac{1}{m} \cdot v^{1/m} \left( \frac{1 - v^n}{1 - v^{1/m}} \right) \nonumber\\[0.3em]
    &= \frac{1 - v^n}{m \left( v^{-1/m} - 1 \right)} = \frac{1 - v^n}{m \left[ (1+i)^{1/m} - 1 \right]}
\end{align}
Berdasarkan definisi suku bunga nominal $i^{(m)} = m[(1+i)^{1/m}-1]$, kita peroleh:
\begin{equation}
    \mathbf{a_{\overline{n}|}^{(m)} = \frac{1 - v^n}{i^{(m)}} = \frac{i}{i^{(m)}} \cdot a_{\overline{n}|i}}
\end{equation}

\subsubsection{Penurunan Nilai Terakumulasi \texorpdfstring{$s_{\overline{n}|}^{(m)}$}{s\_n\textasciicircum(m)}}
\begin{equation}
    \mathbf{s_{\overline{n}|}^{(m)} = a_{\overline{n}|}^{(m)} \cdot (1+i)^n = \frac{(1+i)^n - 1}{i^{(m)}} = \frac{i}{i^{(m)}} \cdot s_{\overline{n}|i}}
\end{equation}

---

\subsection{Anuitas Awal: Notasi \texorpdfstring{$\ddot{a}_{\overline{n}|}^{(m)}$}{a\_ddot\_n\textasciicircum(m)} dan \texorpdfstring{$\ddot{s}_{\overline{n}|}^{(m)}$}{s\_ddot\_n\textasciicircum(m)}}

Misalkan pembayaran sebesar $\frac{1}{m}$ dilakukan pada setiap awal sub-periode $\frac{1}{m}$, yaitu pada waktu:
\begin{equation*}
    t = 0, \frac{1}{m}, \frac{2}{m}, \dots, n - \frac{1}{m}
\end{equation*}

\subsubsection{Penurunan Matematis Nilai Sekarang \texorpdfstring{$\ddot{a}_{\overline{n}|}^{(m)}$}{a\_ddot\_n\textasciicircum(m)}}
\begin{align}
    \ddot{a}_{\overline{n}|}^{(m)} &= \frac{1}{m} \left( 1 + v^{1/m} + v^{2/m} + \dots + v^{n - 1/m} \right) \nonumber\\[0.3em]
    &= \frac{1}{m} \left( \frac{1 - v^n}{1 - v^{1/m}} \right) = \frac{1 - v^n}{m \left( 1 - v^{1/m} \right)}
\end{align}
Berdasarkan definisi tingkat diskonto nominal $d^{(m)} = m[1 - v^{1/m}]$, kita peroleh:
\begin{equation}
    \mathbf{\ddot{a}_{\overline{n}|}^{(m)} = \frac{1 - v^n}{d^{(m)}} = \frac{i}{d^{(m)}} \cdot a_{\overline{n}|i} = \frac{d}{d^{(m)}} \cdot \ddot{a}_{\overline{n}|i}}
\end{equation}

\subsubsection{Penurunan Nilai Terakumulasi \texorpdfstring{$\ddot{s}_{\overline{n}|}^{(m)}$}{s\_ddot\_n\textasciicircum(m)}}
\begin{equation}
    \mathbf{\ddot{s}_{\overline{n}|}^{(m)} = \ddot{a}_{\overline{n}|}^{(m)} \cdot (1+i)^n = \frac{(1+i)^n - 1}{d^{(m)}} = \frac{i}{d^{(m)}} \cdot s_{\overline{n}|i} = \frac{d}{d^{(m)}} \cdot \ddot{s}_{\overline{n}|i}}
\end{equation}

---

\subsection{Perpetuitas yang Dibayar \texorpdfstring{$m$}{m} Kali per Periode}

Dengan mengambil limit $n \to \infty$ ($v^n \to 0$):
\begin{equation}
    a_{\overline{\infty}|}^{(m)} = \lim_{n \to \infty} a_{\overline{n}|}^{(m)} = \mathbf{\frac{1}{i^{(m)}}}
\end{equation}
\begin{equation}
    \ddot{a}_{\overline{\infty}|}^{(m)} = \lim_{n \to \infty} \ddot{a}_{\overline{n}|}^{(m)} = \mathbf{\frac{1}{d^{(m)}}}
\end{equation}

\begin{notebox}[Hubungan Antar-Kuantitas dan Hirarki Bunga]
Ingat kembali relasi ketaksamaan fundamental antara suku bunga nominal, suku bunga efektif, dan force of interest:
\begin{equation}
    d < d^{(m)} < \delta < i^{(m)} < i \quad (\text{untuk } m > 1)
\end{equation}
Oleh karena itu, faktor pengali selalu bernilai lebih besar dari satu:
\begin{equation}
    \frac{i}{d^{(m)}} > \frac{i}{i^{(m)}} > 1 \implies \ddot{a}_{\overline{n}|}^{(m)} > a_{\overline{n}|}^{(m)} > a_{\overline{n}|i}
\end{equation}
Selain itu, hubungan antara versi due dan immediate adalah:
\begin{equation}
    \ddot{a}_{\overline{n}|}^{(m)} = (1+i)^{1/m} \cdot a_{\overline{n}|}^{(m)} \qquad \text{dan} \qquad \ddot{s}_{\overline{n}|}^{(m)} = (1+i)^{1/m} \cdot s_{\overline{n}|}^{(m)}
\end{equation}
\end{notebox}

---

\subsection{Contoh Soal Pembahasan Kasus \texorpdfstring{$m > 1$}{m > 1}}

\subsubsection{Contoh Soal Pinjaman dengan Notasi \texorpdfstring{$a_{\overline{n}|}^{(m)}$}{a\_n\textasciicircum(m)}}

\begin{examplebox}[Contoh 2 Diselesaikan dengan Pendekatan Aljabar (Slide Hal. 17--18)]
Pinjaman $\$3.000$ selama 5 tahun dengan bunga nominal $10\%$ \textit{convertible semiannually}. Angsuran $R$ dibayarkan setiap akhir kuartal. Hitung $R$.
\end{examplebox}

\begin{proof}[\textbf{Penyelesaian Aljabar:}]
\begin{itemize}
    \item Periode konversi bunga dasar: Semesteran.
    \item Jangka waktu dalam satuan semester: $n = 5 \times 2 = 10\text{ semester}$.
    \item Suku bunga efektif per semester: $i = \frac{10\%}{2} = 5\% = 0.05$.
    \item Frekuensi pembayaran per semester: $m = 2$ (karena 1 semester $= 2$ kuartal).
    \item Total pembayaran per semester: Jika angsuran per kuartal $= R$, maka total aliran kas per 1 semester $= 2R$.
\end{itemize}
Persamaan nilai sekarang:
\begin{equation*}
    PV = (2R) \cdot a_{\overline{10}|0.05}^{(2)} = 3000 \implies R = \frac{1500}{a_{\overline{10}|0.05}^{(2)}}
\end{equation*}
Hitung suku bunga nominal per semester yang dikonversi 2 kali per semester ($i^{(2)}$):
\begin{equation*}
    i^{(2)} = 2 \left[ (1 + 0.05)^{1/2} - 1 \right] = 2 [1.024695076 - 1] = 0.04939015
\end{equation*}
Maka:
\begin{align*}
    a_{\overline{10}|0.05}^{(2)} &= \frac{i}{i^{(2)}} \cdot a_{\overline{10}|0.05} = \frac{0.05}{0.04939015} \cdot \left(\frac{1 - (1.05)^{-10}}{0.05}\right) \\
    &= \frac{1 - 0.61391325}{0.04939015} = \frac{0.38608675}{0.04939015} \approx 7.81708
\end{align*}
Maka besar angsuran per kuartal adalah:
\begin{equation*}
    R = \frac{1500}{7.81708} = \mathbf{191.89}
\end{equation*}
\end{proof}

---

\subsubsection{Contoh Soal Suku Bunga Perpetuitas Due \texorpdfstring{$m$}{m}-thly}

\begin{examplebox}[Contoh Perpetuitas (Slide Hal. 19--20)]
Pada tingkat suku bunga efektif tahunan berapakah ($i$), nilai sekarang dari sederetan pembayaran $\$1$ pada setiap 6 bulan selamanya adalah $\$10$, di mana pembayaran pertama dimulai \textbf{segera saat ini} ($t=0$)?
\end{examplebox}

\begin{proof}[\textbf{Metode 1: Menggunakan Notasi Perpetuitas Due $\ddot{a}_{\overline{\infty}|}^{(m)}$}]
\begin{itemize}
    \item Periode konversi bunga dasar: 1 Tahun.
    \item Frekuensi pembayaran per tahun: $m = 2$ kali per tahun.
    \item Total pembayaran per tahun: $2 \times \$1 = \$2$.
    \item Karena pembayaran pertama dimulai segera ($t=0$), ini adalah \textbf{perpetuitas-due}:
\end{itemize}
\begin{equation*}
    PV = 2 \cdot \ddot{a}_{\overline{\infty}|}^{(2)} = 2 \cdot \frac{1}{d^{(2)}} = 10 \implies d^{(2)} = \frac{2}{10} = 0.20
\end{equation*}
Ingat definisi $d^{(2)} = 2 \left[ 1 - (1-d)^{1/2} \right]$:
\begin{equation*}
    2 \left[ 1 - (1-d)^{1/2} \right] = 0.20 \implies 1 - (1-d)^{1/2} = 0.10 \implies (1-d)^{1/2} = 0.90
\end{equation*}
Kuadratkan kedua ruas:
\begin{equation*}
    1 - d = 0.81 \implies d = 0.19
\end{equation*}
Hubungkan tingkat diskonto efektif $d$ dengan suku bunga efektif tahunan $i$:
\begin{equation*}
    i = \frac{d}{1-d} = \frac{0.19}{0.81} = \frac{19}{81} \approx \mathbf{0.2345679} = \mathbf{23.46\%}
\end{equation*}

\vspace{0.5em}
\textbf{Metode 2: Deret Geometri Langsung (Nilai Waktu Uang)}
Biarkan $v = (1+i)^{-1}$ menyatakan faktor diskonto tahunan. Pembayaran $\$1$ terjadi pada $t = 0, 0.5, 1.0, 1.5, \dots$:
\begin{equation*}
    PV = 1 + v^{0.5} + v^1 + v^{1.5} + \dots = \frac{1}{1 - v^{0.5}} = 10
\end{equation*}
\begin{equation*}
    1 - v^{0.5} = \frac{1}{10} = 0.10 \implies v^{0.5} = 0.90 \implies (1+i)^{-0.5} = 0.90
\end{equation*}
\begin{equation*}
    1 + i = (0.90)^{-2} = \frac{1}{0.81} \approx 1.2345679 \implies i = \mathbf{23.46\%}
\end{equation*}
\end{proof}

---

\subsubsection{Contoh Soal Komparasi Beasiswa (\textit{Mahpee Scholarship vs Fellowship})}

\begin{examplebox}[Problem Beasiswa Mahpee (Slide Hal. 21)]
Mahpee ditawari beasiswa (\textit{scholarship}) dengan pembayaran bulanan konstan selama 4 tahun. Pembayaran pertama dilakukan pada 30 September 2000, dan setiap tahun total pembayarannya adalah $\$5.904$. 

Mahpee ingin membandingkan beasiswa ini dengan \textit{fellowship} dari perguruan tinggi lain yang menawarkan pembayaran $\$5.800$ pada 1 September untuk tahun 2000, 2001, 2002, dan 2003.

Asumsikan suku bunga efektif tahunan adalah $i = 4\%$. Bandingkan nilai sekarang kedua tawaran tersebut pada \textbf{1 September 2000}.
\end{examplebox}

\begin{proof}[\textbf{Langkah Analisis \& Perhitungan:}]
Titik evaluasi: \textbf{1 September 2000} ($t = 0$), dengan $i = 0.04$.

\textbf{1. Valuasi Fellowship:}
\begin{itemize}
    \item Pembayaran dilakukan setiap 1 September (awal tahun) selama 4 tahun: 2000, 2001, 2002, 2003.
    \item Ini adalah \textbf{annuity-due} tahunan sebesar $\$5.800$ per tahun selama $n = 4$ tahun.
\end{itemize}
\begin{align*}
    PV_{\text{Fellowship}} &= 5800 \cdot \ddot{a}_{\overline{4}|0.04} = 5800 \cdot (1+i) \cdot a_{\overline{4}|0.04} \\
    &= 5800 \cdot (1.04) \cdot \left( \frac{1 - (1.04)^{-4}}{0.04} \right) \\
    &= 5800 \cdot (1.04) \cdot 3.6298952 = 5800 \cdot 3.775091 = \mathbf{\$21.895,53}
\end{align*}

\textbf{2. Valuasi Scholarship:}
\begin{itemize}
    \item Pembayaran bulanan dimulai 30 September 2000 (akhir bulan ke-1 dari titik evaluasi 1 September 2000).
    \item Ini adalah \textbf{annuity-immediate} yang dibayarkan $m=12$ kali setahun selama $n=4$ tahun, dengan total pembayaran per tahun $\$5.904$ (yaitu $\frac{5904}{12} = \$492$ per bulan).
\end{itemize}
Hitung suku bunga nominal $i^{(12)}$:
\begin{equation*}
    i^{(12)} = 12 \left[ (1.04)^{1/12} - 1 \right] \approx 12 [1.00327374 - 1] = 0.03928488
\end{equation*}
Maka:
\begin{align*}
    PV_{\text{Scholarship}} &= 5904 \cdot a_{\overline{4}|0.04}^{(12)} = 5904 \cdot \frac{1 - (1.04)^{-4}}{i^{(12)}} \\
    &= 5904 \cdot \frac{0.1451958}{0.03928488} = 5904 \cdot 3.695972 = \mathbf{\$21.820,83}
\end{align*}

\textbf{Kesimpulan Analisis Finansial:}
\begin{itemize}
    \item Nilai sekarang Fellowship: \textbf{\$21.895,53}
    \item Nilai sekarang Scholarship: \textbf{\$21.820,83}
    \item \textbf{Keputusan:} Meskipun total uang nominal dari scholarship lebih besar ($4 \times 5904 = \$23.616$) dibandingkan fellowship ($4 \times 5800 = \$23.200$), \textbf{Fellowship memiliki nilai ekonomi yang lebih tinggi} pada 1 September 2000 sebesar $\$74,70$ karena dana fellowship diterima di muka pada awal tahun.
\end{itemize}
\end{proof}

---

\section{Kasus 3 Aljabar: Anuitas Kontinu (\textit{Continuous Annuities})}

\subsection{Motivasi dan Definisi Matematis}
Ketika frekuensi pembayaran meningkat tanpa batas ($m \to \infty$), arus kas tidak lagi terjadi pada titik-titik diskret, melainkan mengalir secara kontinu sepanjang waktu (\textit{continuous cash flow}). 

Konsep anuitas kontinu sangat penting dalam:
\begin{itemize}
    \item Teori aktuaria matematika asuransi jiwa (asuransi kematian kontinu, cadangan premi).
    \item Aproksimasi yang sangat akurat untuk arus kas dengan frekuensi transaksi harian atau per jam.
    \item Analisis teoritis menggunakan kalkulus dan persamaan diferensial.
\end{itemize}

\begin{definitionbox}[Simbol Standar Aktuaria untuk Anuitas Kontinu]
\begin{itemize}[leftmargin=1.5em]
    \item $\bar{a}_{\overline{n}|}$: Nilai sekarang dari anuitas yang dibayarkan secara kontinu dengan laju konstan $1$ per unit waktu selama $n$ periode.
    \item $\bar{s}_{\overline{n}|}$: Nilai akumulasi pada $t = n$ dari anuitas kontinu dengan laju konstan $1$ per unit waktu selama $n$ periode.
\end{itemize}
\end{definitionbox}

---

\subsection{Penurunan Nilai Sekarang (\texorpdfstring{$\bar{a}_{\overline{n}|}$}{a\_bar\_n|})}

Arus kas sebesar $dt$ dibayarkan pada saat $t$. Nilai sekarang dari pembayaran diferensial ini adalah $v^t \, dt$. Dengan mengintegrasikan seluruh arus kas dari $t = 0$ hingga $t = n$:
\begin{equation}
    \bar{a}_{\overline{n}|} = \int_{0}^{n} v^t \, dt
\end{equation}
Ingat bahwa $v^t = (1+i)^{-t} = e^{-t \ln(1+i)} = e^{-\delta t}$, di mana $\delta = \ln(1+i)$ adalah \textit{force of interest}:
\begin{align}
    \bar{a}_{\overline{n}|} &= \int_{0}^{n} e^{-\delta t} \, dt = \left[ -\frac{e^{-\delta t}}{\delta} \right]_{0}^{n} = -\frac{e^{-n\delta}}{\delta} - \left( -\frac{1}{\delta} \right) \nonumber\\[0.4em]
    &= \mathbf{\frac{1 - e^{-n\delta}}{\delta} = \frac{1 - v^n}{\delta}}
\end{align}

\subsubsection{Konsistensi Melalui Pendekatan Limit}
Kita juga dapat membuktikan rumus ini dengan mengambil limit $m \to \infty$ dari anuitas diskret $a_{\overline{n}|}^{(m)}$ dan $\ddot{a}_{\overline{n}|}^{(m)}$:
\begin{equation}
    \bar{a}_{\overline{n}|} = \lim_{m \to \infty} a_{\overline{n}|}^{(m)} = \lim_{m \to \infty} \frac{1 - v^n}{i^{(m)}} = \frac{1 - v^n}{\lim_{m \to \infty} i^{(m)}} = \frac{1 - v^n}{\delta}
\end{equation}
\begin{equation}
    \bar{a}_{\overline{n}|} = \lim_{m \to \infty} \ddot{a}_{\overline{n}|}^{(m)} = \lim_{m \to \infty} \frac{1 - v^n}{d^{(m)}} = \frac{1 - v^n}{\lim_{m \to \infty} d^{(m)}} = \frac{1 - v^n}{\delta}
\end{equation}
Karena $\lim_{m \to \infty} i^{(m)} = \lim_{m \to \infty} d^{(m)} = \delta$.

\begin{definitionbox}[Rumus Hubungan Antara $\bar{a}_{\overline{n}|}$, $a_{\overline{n}|}$, dan $\ddot{a}_{\overline{n}|}$]
\begin{equation}
    \mathbf{\bar{a}_{\overline{n}|} = \frac{i}{\delta} \cdot a_{\overline{n}|i} = \frac{d}{\delta} \cdot \ddot{a}_{\overline{n}|i}}
\end{equation}
\end{definitionbox}

---

\subsection{Penurunan Nilai Terakumulasi (\texorpdfstring{$\bar{s}_{\overline{n}|}$}{s\_bar\_n|})}

Nilai akumulasi pada waktu $t = n$ diperoleh dengan mengintegrasikan arus kas yang telah diakumulasikan:
\begin{align}
    \bar{s}_{\overline{n}|} &= \int_{0}^{n} (1+i)^{n-t} \, dt = \int_{0}^{n} (1+i)^x \, dx = \int_{0}^{n} e^{\delta x} \, dx \nonumber\\[0.4em]
    &= \left[ \frac{e^{\delta x}}{\delta} \right]_{0}^{n} = \mathbf{\frac{e^{n\delta} - 1}{\delta} = \frac{(1+i)^n - 1}{\delta}}
\end{align}
Hubungan dengan nilai sekarang:
\begin{equation}
    \bar{s}_{\overline{n}|} = \bar{a}_{\overline{n}|} \cdot (1+i)^n = \bar{a}_{\overline{n}|} \cdot e^{n\delta}
\end{equation}
\begin{equation}
    \mathbf{\bar{s}_{\overline{n}|} = \frac{i}{\delta} \cdot s_{\overline{n}|i} = \frac{d}{\delta} \cdot \ddot{s}_{\overline{n}|i}}
\end{equation}

---

\subsection{Perpetuitas Kontinu (\texorpdfstring{$\bar{a}_{\overline{\infty}|}$}{a\_bar\_infinity})}

Jika aliran dana kontinu berlangsung selamanya ($n \to \infty$):
\begin{equation}
    \mathbf{\bar{a}_{\overline{\infty}|} = \lim_{n \to \infty} \frac{1 - e^{-n\delta}}{\delta} = \frac{1}{\delta}}
\end{equation}

---

\subsection{Perspektif Kalkulus \& Persamaan Diferensial Biasa (PDB)}

\subsubsection{1. Diferensiasi Terhadap Waktu untuk Nilai Akumulasi}
Misalkan saldo dana tabungan kontinu pada saat $t$ dinotasikan dengan $\bar{s}_{\overline{t}|}$. Berdasarkan Teorema Dasar Kalkulus:
\begin{equation}
    \frac{d}{dt} \bar{s}_{\overline{t}|} = \frac{d}{dt} \int_{0}^{t} (1+i)^x \, dx = (1+i)^t
\end{equation}
Karena $(1+i)^t = 1 + \left((1+i)^t - 1\right) = 1 + \delta \left( \frac{(1+i)^t - 1}{\delta} \right) = 1 + \delta \bar{s}_{\overline{t}|}$, maka:
\begin{equation}
    \mathbf{\frac{d}{dt} \bar{s}_{\overline{t}|} = 1 + \delta \bar{s}_{\overline{t}|}}
\end{equation}

\begin{notebox}[Interpretasi Ekonomi Dinamis:]
Persamaan diferensial di atas menunjukkan bahwa \textbf{laju pertumbuhan saldo tabungan} ($\frac{d}{dt}\bar{s}_{\overline{t}|}$) terdiri dari dua komponen:
\begin{enumerate}
    \item \textbf{Laju aliran dana masuk baru:} Sebesar $1$ per satuan waktu.
    \item \textbf{Laju perolehan bunga berjalan:} Sebesar $\delta \cdot \bar{s}_{\overline{t}|}$ atas saldo yang telah terkumpul.
\end{enumerate}
\end{notebox}

\subsubsection{2. Diferensiasi Terhadap Waktu untuk Nilai Sekarang}
Untuk nilai sekarang anuitas kontinu $\bar{a}_{\overline{t}|}$:
\begin{equation}
    \frac{d}{dt} \bar{a}_{\overline{t}|} = \frac{d}{dt} \int_{0}^{t} v^x \, dx = v^t
\end{equation}
Karena $v^t = 1 - (1 - v^t) = 1 - \delta \left( \frac{1 - v^t}{\delta} \right) = 1 - \delta \bar{a}_{\overline{t}|}$, maka:
\begin{equation}
    \mathbf{\frac{d}{dt} \bar{a}_{\overline{t}|} = 1 - \delta \bar{a}_{\overline{t}|}}
\end{equation}

\subsubsection{3. Rekonstruksi Rumus melalui Solusi Persamaan Diferensial}
Misalkan $y(t) = \bar{a}_{\overline{t}|}$. Maka persamaan diferensialnya adalah:
\begin{equation*}
    y' + \delta y = 1, \quad \text{dengan syarat awal } y(0) = 0
\end{equation*}
\begin{itemize}
    \item Solusi homogen: $y_h' + \delta y_h = 0 \implies y_h(t) = C e^{-\delta t}$.
    \item Solusi partikular: Tebak konstanta $y_p = A \implies 0 + \delta A = 1 \implies A = \frac{1}{\delta}$.
    \item Solusi umum: $y(t) = C e^{-\delta t} + \frac{1}{\delta}$.
    \item Syarat awal $y(0) = 0$: $C (1) + \frac{1}{\delta} = 0 \implies C = -\frac{1}{\delta}$.
\end{itemize}
Sehingga diperoleh:
\begin{equation*}
    y(t) = \bar{a}_{\overline{t}|} = -\frac{1}{\delta} e^{-\delta t} + \frac{1}{\delta} = \mathbf{\frac{1 - e^{-\delta t}}{\delta}} \quad \text{(Q.E.D.)}
\end{equation*}

---

\subsection{Contoh Soal Anuitas Kontinu}

\begin{examplebox}[Contoh Penentuan Force of Interest (Slide Hal. 26)]
Hitung besarnya \textit{force of interest} $\delta$ jika diketahui hubungan:
\begin{equation*}
    \bar{s}_{\overline{20}|} = 3 \cdot \bar{s}_{\overline{10}|}
\end{equation*}
\end{examplebox}

\begin{proof}[\textbf{Langkah Penyelesaian:}]
Substitusikan rumus anuitas kontinu $\bar{s}_{\overline{n}|} = \frac{e^{n\delta} - 1}{\delta}$:
\begin{equation*}
    \frac{e^{20\delta} - 1}{\delta} = 3 \cdot \frac{e^{10\delta} - 1}{\delta}
\end{equation*}
Karena suku bunga $\delta > 0$, kita dapat mengalikan kedua ruas dengan $\delta$:
\begin{equation*}
    e^{20\delta} - 1 = 3 (e^{10\delta} - 1)
\end{equation*}
Faktorkan bentuk selisih kuadrat $e^{20\delta} - 1 = (e^{10\delta} - 1)(e^{10\delta} + 1)$:
\begin{equation*}
    (e^{10\delta} - 1)(e^{10\delta} + 1) = 3 (e^{10\delta} - 1)
\end{equation*}
Pindahkan semua suku ke ruas kiri:
\begin{equation*}
    (e^{10\delta} - 1)\left[ (e^{10\delta} + 1) - 3 \right] = 0
\end{equation*}
\begin{equation*}
    (e^{10\delta} - 1)(e^{10\delta} - 2) = 0
\end{equation*}
Terdapat dua kemungkinan solusi:
\begin{enumerate}
    \item $e^{10\delta} - 1 = 0 \implies e^{10\delta} = 1 \implies \delta = 0$ (Ditolak, karena kasus trivial tanpa bunga).
    \item $e^{10\delta} - 2 = 0 \implies e^{10\delta} = 2$.
\end{enumerate}
Ambil logaritma natural pada kedua ruas:
\begin{equation*}
    10\delta = \ln(2) \implies \delta = \frac{\ln(2)}{10} \approx \frac{0.693147}{10} = \mathbf{0.069315} = \mathbf{6.9315\%}
\end{equation*}
Jadi, \textit{force of interest} yang memenuhi adalah $\mathbf{6.9315\%}$.
\end{proof}

---

\section{Rangkuman Komparatif \& Tabel Formula Lengkap}

Berikut adalah tabel matriks perbandingan komprehensif seluruh rumus anuitas umum:

\begin{table}[htbp]
\centering
\small
\renewcommand{\arraystretch}{1.6}
\adjustbox{max width=\textwidth}{%
\begin{tabular}{lcccc}
\toprule
\textbf{Jenis Kasus Anuitas} & \textbf{Frekuensi Pembayaran} & \textbf{Nilai Sekarang ($PV$)} & \textbf{Nilai Akumulasi ($FV$)} & \textbf{Perpetuitas ($PV_{\infty}$)} \\
\midrule
\textbf{Kasus Standar} & 1 kali per periode bunga & $a_{\overline{n}|i} = \frac{1-v^n}{i}$ & $s_{\overline{n}|i} = \frac{(1+i)^n-1}{i}$ & $\frac{1}{i}$ \\
\midrule
\textbf{Pembayaran Lebih Jarang} & 1 kali per $k$ periode & \multirow{2}{*}{$\dfrac{a_{\overline{n}|i}}{s_{\overline{k}|i}}$} & \multirow{2}{*}{$\dfrac{s_{\overline{n}|i}}{s_{\overline{k}|i}}$} & \multirow{2}{*}{$\dfrac{1}{i \cdot s_{\overline{k}|i}} = \dfrac{1}{(1+i)^k-1}$} \\
(\textit{Immediate}, $k > 1$) & (pada $t = k, 2k, \dots, n$) & & & \\
\midrule
\textbf{Pembayaran Lebih Jarang} & 1 kali per $k$ periode & \multirow{2}{*}{$\dfrac{a_{\overline{n}|i}}{a_{\overline{k}|i}}$} & \multirow{2}{*}{$\dfrac{s_{\overline{n}|i}}{a_{\overline{k}|i}}$} & \multirow{2}{*}{$\dfrac{1}{i \cdot a_{\overline{k}|i}} = \dfrac{1}{1 - v^k}$} \\
(\textit{Due}, $k > 1$) & (pada $t = 0, k, \dots, n-k$) & & & \\
\midrule
\textbf{Pembayaran Lebih Sering} & $m$ kali per periode & \multirow{2}{*}{$a_{\overline{n}|}^{(m)} = \dfrac{1-v^n}{i^{(m)}} = \dfrac{i}{i^{(m)}} a_{\overline{n}|}$} & \multirow{2}{*}{$s_{\overline{n}|}^{(m)} = \dfrac{(1+i)^n-1}{i^{(m)}} = \dfrac{i}{i^{(m)}} s_{\overline{n}|}$} & \multirow{2}{*}{$\dfrac{1}{i^{(m)}}$} \\
(\textit{Immediate}, $m > 1$) & (sebesar $1/m$ di akhir) & & & \\
\midrule
\textbf{Pembayaran Lebih Sering} & $m$ kali per periode & \multirow{2}{*}{$\ddot{a}_{\overline{n}|}^{(m)} = \dfrac{1-v^n}{d^{(m)}} = \dfrac{i}{d^{(m)}} a_{\overline{n}|}$} & \multirow{2}{*}{$\ddot{s}_{\overline{n}|}^{(m)} = \dfrac{(1+i)^n-1}{d^{(m)}} = \dfrac{i}{d^{(m)}} s_{\overline{n}|}$} & \multirow{2}{*}{$\dfrac{1}{d^{(m)}}$} \\
(\textit{Due}, $m > 1$) & (sebesar $1/m$ di awal) & & & \\
\midrule
\textbf{Pembayaran Kontinu} & Kontinu ($m \to \infty$) & \multirow{2}{*}{$\bar{a}_{\overline{n}|} = \dfrac{1-v^n}{\delta} = \dfrac{i}{\delta} a_{\overline{n}|}$} & \multirow{2}{*}{$\bar{s}_{\overline{n}|} = \dfrac{(1+i)^n-1}{\delta} = \dfrac{i}{\delta} s_{\overline{n}|}$} & \multirow{2}{*}{$\dfrac{1}{\delta}$} \\
(\textit{Continuous}) & (laju 1 per periode) & & & \\
\bottomrule
\end{tabular}%
}
\caption{Matriks Komparasi Formula Anuitas Umum dan Hubungan Simbol Aktuaria}
\end{table}

---

\section{Latihan Soal Mandiri Bergradasi \& Pembahasan Lengkap}

\begin{examplebox}[Latihan Soal 1: Nilai Sekarang Anuitas Pembayaran Jarang]
Sebuah warisan memberikan hak pembayaran sebesar $\$5.000$ yang dibayarkan \textbf{setiap 2 tahun sekali} selama 20 tahun (total 10 kali pembayaran), dengan pembayaran pertama dilakukan 2 tahun dari sekarang. Jika suku bunga efektif tahunan adalah $6\%$, hitunglah nilai sekarang ($PV$) dari warisan tersebut.
\end{examplebox}

\begin{proof}[\textbf{Pembahasan:}]
\textbf{Metode Aljabar ($k > 1$):}
\begin{itemize}
    \item Unit waktu: Tahunan, $i = 0.06$.
    \item Interval pembayaran: $k = 2\text{ tahun}$.
    \item Total durasi: $n = 20\text{ tahun}$.
    \item Besar pembayaran per interval: $\$5.000$.
\end{itemize}
\begin{equation*}
    PV = 5000 \cdot \frac{a_{\overline{20}|0.06}}{s_{\overline{2}|0.06}}
\end{equation*}
Hitung masing-masing kuantitas:
\begin{equation*}
    a_{\overline{20}|0.06} = \frac{1 - (1.06)^{-20}}{0.06} = \frac{1 - 0.3118047}{0.06} = 11.469921
\end{equation*}
\begin{equation*}
    s_{\overline{2}|0.06} = \frac{(1.06)^2 - 1}{0.06} = \frac{1.1236 - 1}{0.06} = 2.06
\end{equation*}
Maka:
\begin{equation*}
    PV = 5000 \cdot \frac{11.469921}{2.06} = 5000 \cdot 5.567923 = \mathbf{\$27.839,61}
\end{equation*}

\textbf{Verifikasi Metode Numerik ($j$):}
Suku bunga efektif per 2 tahun: $j = (1.06)^2 - 1 = 0.1236$.
\begin{equation*}
    PV = 5000 \cdot a_{\overline{10}|0.1236} = 5000 \cdot \left( \frac{1 - (1.1236)^{-10}}{0.1236} \right) = 5000 \cdot 5.567923 = \mathbf{\$27.839,61} \quad \checkmark
\end{equation*}
\end{proof}

\vspace{1em}

\begin{examplebox}[Latihan Soal 2: Penentuan Suku Bunga Nominal dari Nilai Sekarang]
Diketahui bahwa nilai sekarang dari anuitas akhir 10 tahun yang membayar $\$100$ pada setiap akhir bulan adalah $\$8.000$. Hitunglah suku bunga nominal tahunan $i^{(12)}$ yang mendasari transaksi ini!
\end{examplebox}

\begin{proof}[\textbf{Pembahasan:}]
\begin{itemize}
    \item Jangka waktu: $n = 10\text{ tahun}$.
    \item Total pembayaran per tahun: $12 \times 100 = \$1.200$.
    \item $PV = 1200 \cdot a_{\overline{10}|}^{(12)} = 8000 \implies a_{\overline{10}|}^{(12)} = \frac{8000}{1200} = \frac{20}{3} \approx 6.6667$.
\end{itemize}
Secara ekuivalen dalam tingkat suku bunga bulanan $j = \frac{i^{(12)}}{12}$ dengan $N = 120$ bulan:
\begin{equation*}
    8000 = 100 \cdot a_{\overline{120}|j} \implies a_{\overline{120}|j} = 80.0
\end{equation*}
Menggunakan aproksimasi numerik (atau kalkulator finansial):
Untuk $j = 0.007$ ($0.7\%$ per bulan): $a_{\overline{120}|0.007} = \frac{1 - (1.007)^{-120}}{0.007} \approx 80.99$.
Untuk $j = 0.0072$ ($0.72\%$ per bulan): $a_{\overline{120}|0.0072} \approx 79.91$.
Melalui interpolasi linier presisi didapat $j \approx 0.007183 = 0.7183\%$ per bulan.
Maka suku bunga nominal tahunan adalah:
\begin{equation*}
    i^{(12)} = 12 \times j = 12 \times 0.7183\% = \mathbf{8.62\%}
\end{equation*}
\end{proof}

\vspace{1em}

\begin{examplebox}[Latihan Soal 3: Persamaan Anuitas Kontinu \& Suku Bunga]
Jika diketahui bahwa $\bar{a}_{\overline{n}|} = 10$ dan $\bar{s}_{\overline{n}|} = 15$, tentukan:
\begin{enumerate}[label=\alph*.]
    \item Besarnya \textit{force of interest} $\delta$.
    \item Lamanya periode $n$.
\end{enumerate}
\end{examplebox}

\begin{proof}[\textbf{Pembahasan:}]
\textbf{Bagian (a): Menentukan $\delta$}
Ingat hubungan resiprokal dasar antara nilai sekarang dan nilai akumulasi:
\begin{equation*}
    \frac{1}{\bar{a}_{\overline{n}|}} - \frac{1}{\bar{s}_{\overline{n}|}} = \frac{\delta}{1 - v^n} - \frac{\delta}{(1+i)^n - 1} = \frac{\delta}{1 - v^n} - \frac{\delta v^n}{1 - v^n} = \frac{\delta (1 - v^n)}{1 - v^n} = \delta
\end{equation*}
Substitusikan nilai yang diketahui:
\begin{equation*}
    \delta = \frac{1}{10} - \frac{1}{15} = \frac{3 - 2}{30} = \mathbf{\frac{1}{30}} \approx \mathbf{0.03333} \quad (3.33\%)
\end{equation*}

\textbf{Bagian (b): Menentukan $n$}
Hubungkan $\bar{s}_{\overline{n}|}$ dan $\bar{a}_{\overline{n}|}$:
\begin{equation*}
    \bar{s}_{\overline{n}|} = \bar{a}_{\overline{n}|} \cdot e^{n\delta} \implies 15 = 10 \cdot e^{n \cdot (1/30)}
\end{equation*}
\begin{equation*}
    e^{n/30} = \frac{15}{10} = 1.5 \implies \frac{n}{30} = \ln(1.5)
\end{equation*}
\begin{equation*}
    n = 30 \cdot \ln(1.5) = 30 \cdot 0.405465 = \mathbf{12.164}\text{ periode}
\end{equation*}
\end{proof}

\end{document}
